\section{问题一：包含色散与反射相移的两束干涉模型}
\label{sec:single-theory}

\subsection{几何关系、单位与适用条件}
将空气、外延层和衬底依次记为介质 $0,1,2$，外延层厚度为 $d$，
入射角 $\theta$ 相对于界面法线定义。计算中以 $\sigma$ 表示题目给定的波数，
单位为 $\mathrm{cm}^{-1}$，并令 $d_c=10^{-4}d$，其中 $d$ 的数值单位为 $\mu\mathrm{m}$。
真空波长满足 $\lambda_{\mu\mathrm m}=10^4/\sigma$。因此，波数与微米厚度相乘时必须显式乘入 $10^{-4}$。
假定受测区域局部平行、材料线性且无磁性；在弱吸收波段，以标量折射率描述有效传播模式。
这一近似不意味着题目已经给定碳化硅的晶型和光轴方向。

\paperfigure{figures/geometry.png}{平行外延层的反射与相位定义；入射角与折射角均相对于法线测量}{0.72}{fig:geometry}

记空气折射率为 $n_0\simeq1$、外延层折射率为 $n(\sigma)$。
Snell 定律及纵向传播因子为
\begin{equation}
 n_0\sin\theta=n(\sigma)\sin\beta,\qquad
 q_\theta(\sigma)=n(\sigma)\cos\beta
 =\sqrt{n^2(\sigma)-n_0^2\sin^2\theta}.
 \label{eq:snell-q}
\end{equation}
两条出射光必须比较同一出射波前上的相位。外延层内部的往返几何光程是
$2nd_c/\cos\beta$，扣除横向位移造成的外部相位差后，有效光程差为
$2nd_c\cos\beta=2d_cq_\theta$，故传播相位差为
\begin{equation}
 \Phi_{\mathrm{prop},\theta}(\sigma)=4\pi d_c\sigma q_\theta(\sigma).
 \label{eq:prop-phase}
\end{equation}
不能直接将内部射线路径长度 $2nd_c/\cos\beta$ 代入反射干涉公式。
上述平面界面处理与经典干涉及薄膜光学框架一致\cite{BornWolf1999,Byrnes2016}。

\subsection{一次反射的复振幅与半波损失}
用 $r_{ij}$、$t_{ij}$ 表示从介质 $i$ 入射到 $j$ 时的电场振幅反射、透射系数，
对每一种偏振分别计算。在无吸收近似下，题目要求保留的两束反射光叠加为
\begin{equation}
 r^{(2)}_\theta=r_{01}+t_{01}t_{10}r_{12}
 \exp\!\left(\mathrm{i}\Phi_{\mathrm{prop},\theta}\right).
 \label{eq:two-beam-amp}
\end{equation}
令 $a=r_{01}$、$b=t_{01}t_{10}r_{12}$，反射率为
\begin{align}
 R^{(2)}_\theta(\sigma)
 &=|a|^2+|b|^2+2|a b|\cos\Phi_\theta(\sigma),\\
 \Phi_\theta(\sigma)
 &=4\pi d_c\sigma q_\theta(\sigma)+\phi_{r,\theta}(\sigma),\qquad
 \phi_{r,\theta}=\arg b-\arg a.
 \label{eq:two-beam-phase}
\end{align}
低折射率向高折射率界面的反射相移以及两界面相移之差，均已包含在
$\phi_{r,\theta}$ 中。若两束反射光的净相移为 $\pi$，则明暗条纹互换，
但相邻同类条纹的级次差仍为 $1$。因而无需知道绝对干涉级次，也不能不加判断地
把所有极大值都设为整数光程级次。吸收界面还可产生随波数变化的连续相移。

对测量背景及条纹振幅作缓变近似，可写成便于估计的形式
\begin{equation}
 R_\theta(\sigma)=B_\theta(\sigma)+A_\theta(\sigma)
 \cos\!\left[4\pi d_cX_\theta(\sigma)+\phi_{r,\theta}(\sigma)\right]
 +\varepsilon_\theta(\sigma),\qquad
 X_\theta(\sigma)=\sigma q_\theta(\sigma).
 \label{eq:dispersive-coordinate}
\end{equation}
这里 $B_\theta$ 和 $A_\theta$ 表示受界面反射及光谱响应共同影响的背景和振幅，
不将其自由变化解释为材料常数的测量。实际算法限制它们的平滑尺度，使其无法吸收条纹本身。

\subsection{从条纹相位得到厚度}
若 $\sigma_a,\sigma_b$ 是相隔 $k$ 个周期的同类相位极值点，且
$\Delta\phi_r=\phi_{r,\theta}(\sigma_b)-\phi_{r,\theta}(\sigma_a)$，则
\begin{equation}
 2d_c\big[X_\theta(\sigma_b)-X_\theta(\sigma_a)\big]
 +\frac{\Delta\phi_r}{2\pi}=k,
 \qquad
 d=\frac{10^4\left(k-\Delta\phi_r/(2\pi)\right)}
 {2\left[X_\theta(\sigma_b)-X_\theta(\sigma_a)\right]}.
 \label{eq:thickness-dispersive}
\end{equation}
相邻极大值或相邻极小值取 $k=1$；相邻极大值与极小值只跨越半周期，取 $k=1/2$。
当 $n$ 及反射相移均近似恒定时，式\eqref{eq:thickness-dispersive}退化为
\begin{equation}
 d=\frac{10^4}{2\,\Delta\sigma\sqrt{n^2-n_0^2\sin^2\theta}}.
 \label{eq:thickness-constant-n}
\end{equation}
这也是厚度、条纹间距和入射角之间的基本换算关系。

色散存在时，局部周期取决于相位对波数的导数，
\begin{align}
 q_{g,\theta}(\sigma)&=\frac{\mathrm{d}X_\theta}{\mathrm{d}\sigma}
 =q_\theta+\frac{\sigma n}{q_\theta}\frac{\mathrm{d}n}{\mathrm{d}\sigma},\\
 \frac{1}{\Delta\sigma}&\simeq2d_cq_{g,\theta}
 +\frac{1}{2\pi}\frac{\mathrm{d}\phi_{r,\theta}}{\mathrm{d}\sigma}.
 \label{eq:group-optical-thickness}
\end{align}
在法向入射时 $q_g=n+\sigma n'=n-\lambda\,\mathrm{d}n/\mathrm{d}\lambda$。
因此，在明显色散的区间以相位折射率 $n$ 直接换算局部条纹周期，会把群光学厚度当成相位光学厚度。
本文使用完整色散坐标拟合相位，避免先平均多个不等间距条纹造成的偏差。

还需区分原始反射率极值与相位极值。由式\eqref{eq:dispersive-coordinate}可得
\begin{equation}
 \frac{\mathrm{d}R}{\mathrm{d}\sigma}
 =B'+A'\cos\Phi-A\Phi'\sin\Phi.
 \label{eq:envelope-peak-shift}
\end{equation}
当背景或包络的导数不可忽略时，原始谱峰不再精确满足 $\sin\Phi=0$。
峰距法适于初值与独立核验；最终估计应联合拟合缓变项和相位，而不能只保留肉眼挑选的若干谱峰。

\subsection{折射率先验及其适用范围}
碳化硅可具有不同晶型；六方晶体的普通光与非常光介电响应不同，
自由载流子和晶格振动又会改变红外复折射率\cite{Lindquist2004}。
Shaffer 的可见区测量揭示了晶型和双折射差异，不能直接据此将可见区公式外推到全部红外波段\cite{Shaffer1971}。
对本文采用的透明窗口，以 Wang 等的普通光色散作为一种明确的条件模型\cite{Wang2013}。
由 RefractiveIndex.INFO 数据记录转录的 $4H$ 普通光关系为\cite{Polyanskiy2024}
\begin{equation}
 n_o^2-1=
 \frac{0.20075\lambda^2}{\lambda^2+12.07224}
 +\frac{5.54861\lambda^2}{\lambda^2-0.02641}
 +\frac{35.65066\lambda^2}{\lambda^2-1268.24708},
 \label{eq:wang-sic}
\end{equation}
其中 $\lambda$ 以 $\mu\mathrm{m}$ 计；该记录声明的范围是 $0.4047$--$5\,\mu\mathrm{m}$。
所以主计算窗口应满足 $\sigma\geq2000\,\mathrm{cm}^{-1}$。
普通光假设提供可复现的参照，并不能据此断定附件样品的实际晶型、取向和掺杂。

硅的参照色散采用 Chandler-Horowitz 与 Amirtharaj 的室温关系\cite{ChandlerHorowitz2005}
\begin{equation}
 n_{\mathrm{Si}}^2(\lambda)=11.67316+\frac{1}{\lambda^2}
 +\frac{0.004482633}{\lambda^2-1.108205^2},
 \label{eq:silicon-index}
\end{equation}
同样以微米为波长单位。原研究覆盖 $450$--$4000\,\mathrm{cm}^{-1}$，
但其折射率精度不能直接转移到掺杂和温度均未给定的附件样品。
数值计算使用的系数、适用范围及核查来源同时保存在支撑材料中。

\subsection{双角观测的作用与可辨识性边界}
同一样品的两个角度共享厚度，却具有不同的 $q_\theta$，因此能约束厚度并检验模型一致性。
不过，这不意味着两条强度谱一定能稳定分离任意折射率函数与厚度。
以常折射率、无反射相位色散的理想情形说明：定义
$p_\theta=1/(2\Delta\sigma_\theta)=d_c\sqrt{n^2-n_0^2\sin^2\theta}$，则
\begin{equation}
 d_c^2=\frac{p_{10}^2-p_{15}^2}
 {n_0^2\big(\sin^2 15^\circ-\sin^2 10^\circ\big)}.
 \label{eq:two-angle-identifiability}
\end{equation}
数学上可以解出 $d_c$，但分子是两个接近数的差。
例如取示意折射率 $n=2.5$ 时，两个角度的 $p$ 相对差仅约 $0.3\%$；
取 $n=3.42$ 时仅约 $0.16\%$。角度标定、反射相位和背景误差都可能超过这一微小差异。
未知色散、消光系数、偏振混合和衬底响应进一步增加了相关性。
实际可辨识性应由参数扫描或剖面目标函数检验，不能只由优化器返回有限协方差判定\cite{Raue2009}。

对于式\eqref{eq:thickness-constant-n}，在 $n_0=1$、小扰动条件下有
\begin{equation}
 \frac{\delta d}{d}\simeq-\frac{\delta(\Delta\sigma)}{\Delta\sigma}
 -\frac{n\,\delta n}{n^2-\sin^2\theta}
 +\frac{\sin\theta\cos\theta}{n^2-\sin^2\theta}\,\delta\theta,
 \label{eq:relative-sensitivity}
\end{equation}
其中 $\delta\theta$ 以弧度计。近法向入射时，折射率的相对变化几乎等比例、反方向地传递到厚度。
因此，本文区分固定光学常数后的拟合不确定性、波段与模型选择的稳定性，以及未知材料常数造成的系统敏感性。
缺少仪器分辨率、重复测量和独立折射率标定时，前两者不构成完整的绝对测量不确定度\cite{JCGM2008}。
